%% ============================================================
%% SECTION 1: INTRODUCTION
%% EarthVision-LM Survey -- Artificial Intelligence Review
%% Template: sn-jnl (Springer Nature) | sn-mathphys-num
%% Written last -- after all content sections are complete
%% ============================================================

\section{Introduction}\label{sec1}

The past decade has witnessed two convergent revolutions in artificial intelligence, underpinned by large-scale datasets~\cite{geoprivacy2024,long2021creating} and novel architectures.
The first is the emergence of deep learning as the dominant paradigm for visual
perception~\cite{lecun2015deep}: convolutional neural networks, and subsequently
vision transformers~\cite{dosovitskiy2021vit}, have achieved superhuman performance
on image classification, object detection~\cite{ren2015faster,cheng2016rotated}, and semantic
segmentation benchmarks. The second is the rise of large language models
(LLMs)~\cite{brown2020gpt3,touvron2023llama}: transformer-based architectures
pre-trained on trillions of tokens have acquired broad world knowledge, instruction
following, and multi-step reasoning capabilities that were previously inaccessible
to automated systems.

The convergence of these two revolutions has produced \emph{multimodal large language
models} (MLLMs)---systems that jointly process visual and linguistic inputs to enable
natural language dialogue grounded in image content. Models such as
Flamingo~\cite{flamingo2022}, BLIP-2~\cite{blip22023}, LLaVA~\cite{llava2023}, LLaVA-1.5~\cite{llava15_2023},
GPT-4V~\cite{gpt4v2023}, Gemini~\cite{gemini2023}, Qwen-VL~\cite{qwen_vl2023}, and Pixtral~\cite{pixtral2024} have demonstrated that a single
model can describe, question, and reason about photographs, charts, and documents
through open-ended text interaction. These systems have transformed the interface
between human users and visual information.

\textbf{Remote sensing as the next frontier.}
Remote sensing (RS) imagery---acquired by satellites, aircraft, and unmanned aerial
vehicles across the full electromagnetic spectrum---provides the primary data
source for global monitoring of Earth's land surface, oceans, atmosphere, and
built environment. The scale and frequency of RS data production has grown
dramatically: the Sentinel constellation~\cite{esa2023sentinel} alone generates
more than 12 TB of new imagery per day, and commercial VHR providers collectively
image the entire Earth's land surface at sub-metre resolution multiple times per
year~\cite{globalsat2024}. Extracting actionable intelligence from this data
torrent requires automated analysis systems capable of semantic interpretation
at a level that exceeds what specialised task-specific models can provide.

The emergence of RS-specific MLLMs---systems that apply the conversational
intelligence of general MLLMs to RS imagery---offers a transformative opportunity:
to replace the pipeline of specialised models (classifier for land cover, detector
for objects, captioner for scenes) with a single conversational system that can
answer any spatially grounded question about any RS image in natural language. Since
the first dedicated RS-MLLM (RSGPT~\cite{rsgpt2023}) appeared in mid-2023, the
field has grown to 42 published papers in 2025 alone, representing a 75\,\% annual
growth rate from 2024 (Sect.~\ref{sec8:pubtrend}).

\textbf{The fundamental tension.}
Despite this explosive growth, a fundamental tension characterises the current
RS-MLLM landscape: the architectures, training paradigms, and evaluation protocols
of general MLLMs are poorly suited to the specific constraints of remote sensing.
RS images differ from internet photographs along every dimension that vision models
exploit---viewpoint (nadir versus ground-level), scale (sub-metre to kilometre
resolution within a single dataset), spectral content (optical, SAR, and
hyperspectral), object orientation (arbitrary rotation), and annotation economics
(expert annotation versus crowdsourcing). A CLIP-trained encoder that describes a
ship accurately from a ground-level photograph fails systematically when the same
ship is presented from 500\,m altitude in an RS tile. An LLM fine-tuned on
internet-scale VQA has no framework for interpreting SAR backscatter or
hyperspectral spectral signatures.

The result, as we demonstrate in this survey, is that RS-MLLMs in 2026 remain
fundamentally \emph{perception-limited} despite encouraging recent progress.
More than 85\,\% process only optical imagery despite RS being
inherently multi-sensor; although recent systems such as
Earth-OneVision~\cite{earthonevision2026} substantially expand sensor and task
coverage (optical, SAR, infrared, multispectral, temporal, and video), current
RS-MLLMs as a field still lack broad hyperspectral (HSI) coverage,
physics-specific modality encoding, systematic cross-sensor evaluation protocols,
and operationally grounded reasoning. OBB output remains substantially underrepresented: although GeoChat, EarthGPT,
and GeoGround demonstrate OBB-capable grounding or detection, this capability is
limited to optical imagery and uncommon across the field; cross-sensor OBB
reasoning is essentially unaddressed; and all 18 models
evaluated in HM-Bench~\cite{hmbench2026} exhibit systematic spectral hallucination when
presented with non-RGB sensor data---a finding consistent with structural encoder limitations
across the evaluated sample. These are not incidental engineering gaps---they reflect
structural mismatches between the RS domain and the general MLLM paradigm.

\textbf{Research questions.}
This survey is organised around six research questions (RQs) that map
directly to its analytical sections:

\begin{itemize}
\item \textbf{RQ1} (Sect.~\ref{sec3}): How have RS-MLLM architectures evolved across vision encoder adaptation, connector design, LLM backbone strategy, sensor modality scope, and spatial output capability?
\item \textbf{RQ2} (Sect.~\ref{sec4}): Which RS tasks and sensor modalities are adequately represented, and which remain underdeveloped?
\item \textbf{RQ3} (Sect.~\ref{sec5},~\ref{sec6}): What architectural and training strategies govern the trade-off between RS-specific adaptation, computational cost, and generalisation?
\item \textbf{RQ4} (Sect.~\ref{sec7}): What hallucination mechanisms are reported or empirically supported in RS-MLLMs, and how do they differ from general MLLM hallucination?
\item \textbf{RQ5} (Sect.~\ref{sec8}): How adequate are current datasets and benchmarks for evaluating spatial, multimodal, and operational RS-MLLM intelligence?
\item \textbf{RQ6} (Sect.~\ref{sec9}): What research directions are necessary to advance RS-MLLMs from descriptive perception toward reliable geospatial intelligence?
\end{itemize}

\textbf{Contributions of this survey.}
EarthVision-LM provides an architecture-first systematic synthesis that connects
RS domain physics, MLLM design, multimodal coverage, hallucination mechanisms,
benchmark evidence, and operational research priorities within a unified analytical
framework (RQ1--RQ6). We make four primary contributions:

\begin{enumerate}

\item \textbf{Five-dimensional RS-MLLM design space} (Sect.~\ref{sec3}). We
introduce a taxonomy organising RS-MLLM architectures across five dimensions:
vision encoder adaptation ($\mathcal{E}$), vision-language connector design
($\mathcal{C}$), LLM backbone and freezing strategy ($\mathcal{L}$), sensor modality
scope ($\mathcal{M}$), and spatial output capability ($\mathcal{S}$). To the best of
our knowledge, no existing RS-MLLM survey maps these five domain challenges to
design dimensions within a single analytical framework, and the taxonomy exposes
systematic design gaps responsible for current RS-MLLM limitations.

\item \textbf{Unified multi-modal coverage analysis} (Sect.~\ref{sec5}). We present
a systematic analysis of RS-MLLM capability across all sensor modalities.
Although recent models such as EarthGPT~\cite{earthgpt2025},
EarthGPT-X~\cite{earthgptx2026}, and Earth-OneVision~\cite{earthonevision2026}
substantially expand sensor coverage, broad HSI capability, physics-grounded
cross-sensor reasoning, and robust spatial reasoning remain underdeveloped
across the field (84.6\,\% of Corpus~B models are optical-only).
Physics-informed analysis explains why these gaps persist and what architectural
changes are required to close them.

\item \textbf{Three-type RS-MLLM hallucination taxonomy} (Sect.~\ref{sec7}). We
introduce a systematic hallucination taxonomy for RS-MLLMs, distinguishing
Type~I (Localization), Type~II (Discrimination), and Type~III (Spectral Hallucination).
To the best of our knowledge, prior RS-MLLM hallucination studies have not
explicitly isolated spectral hallucination as a distinct failure category;
our taxonomy extends prior localisation/discrimination frameworks by
incorporating spectral failure evidence from HM-Bench~\cite{hmbench2026},
which evaluates 18 MLLMs on HSI-specific tasks and finds systematic failure
across all 18 evaluated models.

\item \textbf{Comprehensive task, dataset, and benchmark quantitative synthesis}
(Sects.~\ref{sec4},~\ref{sec8}). We systematically analyse task coverage across 26
models and 9 task types, catalogue 13 training datasets and 8 evaluation benchmarks,
and present a structured quantitative synthesis grounded in 97 primary
RS-MLLM papers, establishing publication trends, venue distribution,
backbone usage, and modality distribution.

\end{enumerate}

\textbf{Scope and methodology.}
Scene understanding surveys~\cite{earthvision2024,hrben2022} provide additional RS context. This survey is PRISMA 2020-compliant~\cite{prisma2020}, with a full systematic
search process documented in Sect.~\ref{sec2:prisma}. The search covered IEEE
Xplore, arXiv (cs.CV, eess.IV), Springer Link, CVPR, ICCV, NeurIPS, AAAI, and
ISPRS proceedings between January 2022 and August 2026. The database search
yielded 847 records; snowballing added 18 new records; after eligibility review
and post-eligibility exclusion, the full reference corpus comprises 210 papers.
Of these, \textbf{97 constitute the primary RS-MLLM synthesis corpus
(Corpus~A)}: 31 model studies, 18 benchmark studies, 14 dataset studies,
9 hallucination studies, and 25 foundation-model papers published within
the 2022--2026 window. The remaining \textbf{113 papers serve as contextual
supporting literature} (RS domain baselines, pre-2022 foundation models)
and are cited throughout but do not enter the primary evidence synthesis.
To prevent denominator confusion, we define three analytical levels (Sect.~\ref{sec2:corpus_def}):
\textbf{Corpus~A} ($n=97$): the primary RS-MLLM synthesis corpus;
\textbf{Corpus~B} ($n=26$ RS-MLLMs): the architecture registry of 26 distinct
models arising from the 31 model papers (5 papers describe variants of
already-counted architectures); \textbf{Corpus~C} ($n=18$ models): the
task-comparable subset with sufficient benchmark data for cross-model
quantitative analysis.
The 97 primary papers span five categories (Table~\ref{tab:corpus_taxonomy});
the 113 contextual papers cover RS object detection, segmentation, SAR, HSI,
change detection, and scene classification providing domain baseline methods.

\textbf{How this survey differs from prior work.}
Table~\ref{tab:survey_comparison} positions EarthVision-LM against the closest existing
RS-MLLM and geospatial-LLM reviews. A recent systematic review in this journal
(Dorabantu \& Badea~\cite{dorabantu2026geollm}, AIR, February 2026) addresses geospatial
reasoning and awareness in LLMs broadly---covering GeoLLMs, GIS integration, spatial
reasoning, and multimodal geospatial systems. EarthVision-LM is complementary and distinct:
its scope is specifically RS-MLLMs, not GeoLLMs in general, and its analytical frame
is architecture-first rather than application-first.

Unlike broader reviews of geospatial large language models and recent RS-MLLM surveys
organised primarily around tasks, training paradigms, or general model evolution,
EarthVision-LM adopts an architecture-first perspective that jointly connects encoder
adaptation, connector design, LLM adaptation, sensor modality scope, and spatial output
capability. This framework is further integrated with quantitative literature synthesis,
cross-modal coverage analysis, a formal hallucination taxonomy, and an operational
research roadmap. Our contribution is not simply additional coverage; it is the
integration of these analytical dimensions into a single
architecture-to-evidence-to-gap framework.

\begin{table*}[t]
\caption{Comparison of EarthVision-LM with the closest existing RS-MLLM and geospatial-LLM
reviews across nine dimensions.
\checkmark\,=\,covered; \checkmark\checkmark\,=\,primary contribution / in-depth coverage;
$\circ$\,=\,partially covered; --\,=\,not covered.
\textbf{Domain}: RS-VLM\,=\,RS vision-language; RS-MLLM\,=\,RS multimodal LLM;
GeoLLM\,=\,geospatial LLM (broader, includes GIS and GeoQA).
\textbf{Sensor physics}: SAR/HSI physics grounded analysis.
\textbf{PRISMA}: explicit PRISMA-2020-compliant systematic review.
}\label{tab:survey_comparison}
\small\setlength\tabcolsep{2.5pt}
\begin{tabular*}{\textwidth}{@{\extracolsep\fill}p{2.8cm}p{1.0cm}p{1.1cm}p{1.1cm}p{1.1cm}p{1.0cm}p{1.0cm}p{1.1cm}p{1.1cm}p{0.85cm}}
\toprule
\textbf{Survey} & \textbf{Domain} &
\textbf{Archit.} &
\textbf{Sensor physics} &
\textbf{Spatial output} &
\textbf{Hallucin.} &
\textbf{Benchmark} &
\textbf{CV-MLLM comp.} &
\textbf{Roadmap} &
\textbf{PRISMA} \\
\midrule
Li~et~al.~\cite{li2024survey}
  & RS-VLM & \checkmark & -- & \checkmark & -- & \checkmark & -- & -- & -- \\
Weng~et~al.~\cite{weng2025isprs}
  & RS-MLLM & \checkmark & -- & \checkmark & -- & \checkmark & -- & -- & -- \\
Dorabantu \& Badea~\cite{dorabantu2026geollm}
  & GeoLLM & -- & -- & -- & -- & -- & -- & \checkmark & -- \\
Ma~et~al.~\cite{ma2026dsorsp}
  & RS-MLLM & \checkmark & -- & \checkmark & -- & \checkmark & \checkmark & \checkmark & -- \\
\textbf{EarthVision-LM (ours)}
  & \textbf{RS-MLLM} & \checkmark\checkmark & \checkmark\checkmark & \checkmark\checkmark & \checkmark\checkmark & \checkmark\checkmark & \checkmark & \checkmark\checkmark & \checkmark \\
\botrule
\end{tabular*}
\footnotetext{Dorabantu \& Badea~\cite{dorabantu2026geollm} (AIR, February 2026) is a systematic review in
\emph{this journal} covering geospatial reasoning and awareness in LLMs broadly---including
GeoLLMs, GIS integration, spatial reasoning, and multimodal geospatial systems---but does not
focus on RS-specific architectures, sensor physics, or hallucination in multi-sensor RS-MLLMs.
Ma~et~al.~\cite{ma2026dsorsp} (arXiv:2607.20284, 22~July~2026) covers RS-MLLM evolution,
multimodal learning, training data, benchmark comparison, and RS-vs-CV-MLLM diagnostic evaluation.
EarthVision-LM's contribution is the \emph{integration} of architectural taxonomy, sensor-physics
analysis, hallucination taxonomy, and evidence-graded synthesis into a single architecture-to-gap
framework---a combination not found in any of the above surveys.}
\end{table*}

\textbf{Claim--evidence discipline.}
Every quantitative claim in this survey is grounded in a specific evidence
source. Key examples: the claim that ``only one RS-MLLM supports conversational
OBB output'' is supported by Table~\ref{tab:master_comparison} and three cited
primary model papers; the claim that ``84.6\,\% of models process optical only''
is derived from the 210-paper coded corpus (Table~\ref{tab:coding_schema}); the
claim that ``all 18 evaluated models exhibit spectral hallucination'' is grounded
in HM-Bench~\cite{hmbench2026}. Evidence tiers (A--D, defined in
Sect.~\ref{sec2:prisma}) are stated throughout. Claims not directly supported by
a cited paper or the corpus analysis are explicitly flagged as author inference
(``consistent with'', ``suggests'') rather than empirical finding.

\textbf{Paper organisation.}
Section~\ref{sec2} establishes the six RS domain challenges and MLLM foundations
that motivate the survey structure. Section~\ref{sec3} addresses \textbf{RQ1} through
the five-dimensional design space and master architecture comparison. Section~\ref{sec4}
addresses \textbf{RQ2} through task-coverage analysis. Sections~\ref{sec5}
and~\ref{sec6} address \textbf{RQ3} through multimodal coverage and PEFT strategies.
Section~\ref{sec7} addresses \textbf{RQ4} through the three-type hallucination taxonomy.
Section~\ref{sec8} addresses \textbf{RQ5} through dataset and benchmark cataloguing
with structured quantitative synthesis. Section~\ref{sec9} addresses \textbf{RQ6}
through eight concrete future research directions. Section~\ref{sec10} concludes.
